\documentclass[11pt,a4paper]{article}
\usepackage[T1]{fontenc}
\usepackage{lmodern}
\usepackage{amsmath,amssymb,amsthm,mathtools,mathrsfs}
\usepackage[margin=25mm]{geometry}
\usepackage{microtype,booktabs,longtable,array,enumitem,hyperref,xurl,graphicx}
\hypersetup{colorlinks=true,linkcolor=blue,citecolor=blue,urlcolor=blue,
 pdftitle={LM4: Finite group actions, loop structures, and profinite extensions},
 pdfauthor={AI-produced by GPT-6 Astra (OpenAI, ChatGPT Pro mode)}}
\setlength{\emergencystretch}{1.5em}
\setlength{\parskip}{3pt}
\newcommand{\Z}{\mathbb Z}\newcommand{\Q}{\mathbb Q}
\newcommand{\R}{\mathbb R}\newcommand{\C}{\mathbb C}\newcommand{\F}{\mathbb F}
\DeclareMathOperator{\Aut}{Aut}\DeclareMathOperator{\Out}{Out}
\DeclareMathOperator{\Inn}{Inn}\DeclareMathOperator{\Hom}{Hom}
\DeclareMathOperator{\id}{id}\DeclareMathOperator{\im}{im}
\DeclareMathOperator{\coker}{coker}\DeclareMathOperator{\Spec}{Spec}
\DeclareMathOperator{\Gal}{Gal}\DeclareMathOperator{\Map}{Map}
\DeclareMathOperator{\Homeo}{Homeo}\DeclareMathOperator{\Diff}{Diff}
\DeclareMathOperator{\rank}{rank}\DeclareMathOperator{\Sec}{Sec}
\DeclareMathOperator{\GL}{GL}\DeclareMathOperator{\PGL}{PGL}
\newenvironment{claim}[3]{\par\medskip\phantomsection\label{cl:#1}%
 \noindent\textbf{#1. [#2] #3.}\ \ignorespaces}{\par\medskip}
\newcommand{\passage}[1]{\par\noindent{\footnotesize\emph{Manuscript:} #1.}\par}
\begin{document}
\paragraph{Unnumbered continuation: the displacement graph and exact compact-family estimates.}
Retain all the actual data of C.81 and C.92: the original closed disk
$\overline D$, its interior $O$, the compatible metric $d$ on
$\overline D$, the closed discrete free proper operator group $Q$,
the boundary-fixed centralizer $\mathscr C_Q$ with a transitive actual
subgroup, and the specified covering $\pi:O\to N=O/Q$.
Fix $x\in O$ and $p=\pi x$. Write $\mathscr F=\mathscr C_Q(x)$.
The quotient surface may be noncompact.

Use the exact pointed lifts, image and inverse maps proved in C.92:
\[
\begin{gathered}
 h_b=\lambda(b),\qquad \pi h_b=b\pi,\qquad h_b(x)=x,
 \qquad h_bq=qh_b\quad(q\in Q),\\
 \delta_b(\pi z)=d_Q(h_bz,z),\qquad
 \mathscr B_p=\{b\in\mathscr I_p:\delta_b\in C_0(N)\},\\
 D:\mathscr F\longrightarrow\mathscr B_p,\qquad
 L:\mathscr B_p\longrightarrow\mathscr F,\qquad
 LD=\id_{\mathscr F},\quad DL=\id_{\mathscr B_p}.
\end{gathered}                                                  \tag{CFT.1}
\]
Here $\mathscr I_p$ is the pointed identity path component in the
compact-open topology of maps and inverses, $D(a)(\pi z)=\pi(az)$,
and $L(b)$ equals $h_b$ on $O$ and the identity on the original
circle. The metric $d_Q(z,w)=\sup_{q\in Q}d(qz,qw)$ retains every
original operator. C.81 proves that it gives the original topology
on $O$. C.92 proves that $b\mapsto h_b$ and $b\mapsto h_b^{-1}$
are continuous for compact-open convergence. In particular each
lift in this note is the specified lift fixing $x$.

Put
\[
 \Delta(b,c)=\sup_{z\in O}d_Q(h_bz,h_cz),\qquad
 d^\uparrow(b,c)=\Delta(b,c)+\Delta(b^{-1},c^{-1}).
                                                               \tag{CFT.2}
\]
By (C.92.15)--(C.92.16), $D$ and $L$ are isometries for
$d^\uparrow$ and the original sum of the uniform map and inverse
distances on the closed disk. The results below give an exact
graph description of this topology and quantitative estimates
on its compact families.

\emph{An explicit metric for the quotient map-and-inverse topology.}
Choose the based plane coordinate $\psi:O\to\mathbb R^2$ from
C.92, and set
$K_n=\pi\psi^{-1}(\{v:|v|\leq n\})$ for $n\geq1$.
Each $K_n$ is compact,
$K_n\subset\operatorname{int}K_{n+1}$, and the interiors cover $N$.
Indeed $\pi$ is open and the corresponding closed balls lie
inside the next open balls. Every compact subset of $N$ lies
in some $K_n$, by a finite subcover of these increasing open
sets. Retain the quotient metric
$\rho(\pi z,\pi w)=\inf_{q\in Q}d_Q(z,qw)$ from C.81.
For pointed homeomorphisms put
\[
\begin{split}
 m(b,c)=\sum_{n=1}^{\infty}2^{-n}\bigg[
 &\min\!\left(1,\sup_{y\in K_n}\rho(by,cy)\right)\\
 &+\min\!\left(1,\sup_{y\in K_n}\rho(b^{-1}y,c^{-1}y)\right)
 \bigg].
\end{split}                                                    \tag{CFT.3}
\]
This is a metric: the triangle inequality follows from that of
$\rho$ and $\min(1,a+b)\leq\min(1,a)+\min(1,b)$; separation
uses $\bigcup K_n=N$. Convergence for $m$ is precisely uniform
convergence of maps and inverses on each $K_n$. In one direction
each nonnegative summand tends to zero. In the other, a finite
initial sum tends to zero and the remaining sum is bounded by
$2\sum_{n>M}2^{-n}$.

For completeness, uniform convergence on compact sets gives
compact-open convergence: for a compact $K$ and an open set
$V$ containing $b(K)$, the compact set $b(K)$ has positive
$\rho$-distance from the closed set $N\setminus V$ whenever
that complement is nonempty. A sufficiently small uniform
change on $K$ therefore stays in $V$.
Conversely fix compact $K$, a map $b$, and $\epsilon>0$.
Cover $K$ by finitely many relatively compact neighborhoods
$U_i$ whose compact closures satisfy
$\rho(bz,by_i)<\epsilon/3$ on $\overline U_i$, for some
$y_i\in K$. The finite compact-open requirements
$c(\overline U_i)\subset B_\rho(by_i,2\epsilon/3)$ imply
$\rho(cy,by)<\epsilon$ for every $y\in K$.
These neighborhoods exist by continuity of $b$ and local
compactness of $N$. Apply the same arguments to inverses.
Thus (CFT.3) gives exactly the stated quotient topology.

\emph{The finite-core estimate with the actual lifts.}
For a nonempty compact $K\subset N$, choose a compact $C\subset O$
with $\pi C=K$. Such a set is obtained from finitely many
local covering sections: shrink their domains to cover $K$
by open sets whose closures lie in the section domains, and
take the section images of the corresponding compact pieces
of $K$. Then $\pi^{-1}K=QC$. Equivariance and $Q$-invariance
give the exact equality
\[
 \sup_{\pi z\in K}d_Q(h_bz,h_cz)
       =\max_{z\in C}d_Q(h_bz,h_cz).
                                                               \tag{CFT.4}
\]
For a point outside that core, the triangle inequality through
the same point $z$ gives
$d_Q(h_bz,h_cz)\leq\delta_b(\pi z)+\delta_c(\pi z)$.
Consequently
\[
 \Delta(b,c)\leq
 \max\left\{
   \max_{z\in C}d_Q(h_bz,h_cz),\quad
   \sup_{y\notin K}\bigl(\delta_b(y)+\delta_c(y)\bigr)
 \right\}.
                                                               \tag{CFT.5}
\]
An empty tail has supremum zero. The core equality uses the
literal equation $h_b(qz)=qh_b(z)$; there is no choice of a
different deck representative on any part of this estimate.

Here is the simultaneous inverse control. Let $P$ be a nonempty compact
metric space and $u\mapsto b_u$ a continuous family in
$\mathscr B_p$ for the quotient topology (CFT.3). For a nonempty compact
$K\subset N$ define its actual image core
\[
              K^P=\{b_u(y):u\in P,\ y\in K\}.
                                                               \tag{CFT.6}
\]
The joint evaluation map is continuous, so $K^P$ is compact.
The exact inverse identity from C.88 is
\[
              \delta_{b_u^{-1}}(y)
                    =\delta_{b_u}(b_u^{-1}y).
                                                               \tag{CFT.7}
\]
For $y\notin K^P$ one has $b_u^{-1}y\notin K$ for every $u$.
Thus a common bound $\delta_{b_u}<\epsilon$ off $K$ gives the
same bound $\delta_{b_u^{-1}}<\epsilon$ off $K^P$.
Choose compact $C^P\subset O$ projecting onto $K^P$ by the
same local-section construction. For all $u,v\in P$, (CFT.5)
applied to maps and inverses yields
\[
\begin{split}
 d^\uparrow(b_u,b_v)\leq
 &\max\left\{\max_{z\in C}d_Q(h_{b_u}z,h_{b_v}z),\ 2\epsilon\right\}\\
 {}+&\max\left\{\max_{z\in C^P}
            d_Q(h_{b_u}^{-1}z,h_{b_v}^{-1}z),\ 2\epsilon\right\}.
\end{split}                                                    \tag{CFT.8}
\]
The inverse core in this formula is exactly (CFT.6); it does
not require a common compact support for the family.

\emph{The exact displacement graph, with both inverse maps.}
Give $C_0(N)$ its supremum-norm topology and define
\[
\begin{gathered}
 \Gamma:\mathscr B_p\longrightarrow\mathscr I_p\times C_0(N),
                  \qquad \Gamma(b)=(b,\delta_b),\\
 \mathscr G_p=\{(b,\delta_b):b\in\mathscr B_p\},\qquad
 \Theta:\mathscr F\longrightarrow\mathscr G_p,
                  \qquad \Theta(a)=\Gamma(D(a)),\\
 \Pi:\mathscr G_p\longrightarrow\mathscr F,
                  \qquad \Pi(b,f)=L(b).
\end{gathered}                                                  \tag{CFT.9}
\]
On the displayed graph $f$ is exactly $\delta_b$. Substitution
in (CFT.1) proves $\Pi\Theta=\id_{\mathscr F}$ and
$\Theta\Pi=\id_{\mathscr G_p}$. The inverse of $\Gamma$ onto
its image is the first-coordinate projection.

These are homeomorphisms for the original topology on
$\mathscr F$ and the product-subspace topology on $\mathscr G_p$.
Here are both continuity proofs. The triangle inequality for
$d_Q$ gives
\[
 \|\delta_b-\delta_c\|_\infty\leq\Delta(b,c),\qquad
 m(b,c)\leq d^\uparrow(b,c),\qquad
 m(b,c)+\|\delta_b-\delta_c\|_\infty
                       \leq2d^\uparrow(b,c).
                                                               \tag{CFT.10}
\]
The middle inequality uses
$\rho(b\pi z,c\pi z)\leq d_Q(h_bz,h_cz)$ and its inverse
counterpart, term by term in (CFT.3), with $\sum2^{-n}=1$.
Together with the C.92 isometry, (CFT.10) proves continuity
of $\Theta$ and $\Gamma$.

To prove the inverse continuity, take a convergent sequence
$(b_n,\delta_{b_n})\to(b,\delta_b)$ in the graph product.
The first coordinates converge for (CFT.3) and the second
coordinates converge in supremum norm. Put
$P=\{0\}\cup\{1/n:n\geq1\}$, with $b_0=b$ and
$b_{1/n}=b_n$. This is a compact continuous quotient family.
The set of its displacement functions is compact in $C_0(N)$,
because it is a norm-convergent sequence together with its limit.
Every compact set of functions in $C_0(N)$ has a common tail:
choose finitely many functions within $\epsilon/2$ of every
function in the set, choose a compact tail core for each of
those finitely many functions with bound $\epsilon/2$, and
take their union together with $\{p\}$. Thus
$\delta_{b_u}<\epsilon$ outside one nonempty compact $K$
for every $u\in P$.

Use $K^P,C,C^P$ from (CFT.6)--(CFT.8).
The joint interior continuity of the lifts from C.92 implies
that both core maxima in (CFT.8), with $u=1/n,v=0$, tend to
zero. To see this directly, the functions
$(u,z)\mapsto d_Q(h_{b_u}z,h_bz)$ on $P\times C$ and
$(u,z)\mapsto d_Q(h_{b_u}^{-1}z,h_b^{-1}z)$ on $P\times C^P$
are continuous and vanish at every $(0,z)$; a finite cover
of each compact core makes the convergence uniform there.
It follows that
$\limsup_n d^\uparrow(b_n,b)\leq4\epsilon$.
As $\epsilon>0$ is arbitrary, this distance tends to zero.
Both the product-subspace topology and $d^\uparrow$ are
metrizable, so sequential continuity proves continuity.
This establishes the claimed homeomorphisms and their
specific inverses on the exact graph (CFT.9).

The graph $\mathscr G_p$ is closed in
$\mathscr I_p\times C_0(N)$. Indeed let
$(b_n,\delta_{b_n})\to(b,f)$ in this product. For each
$y=\pi z$, compact-open lift continuity and the compatibility
of $d_Q$ give
\[
  \delta_{b_n}(y)=d_Q(h_{b_n}z,z)
                \longrightarrow d_Q(h_bz,z)=\delta_b(y).
                                                               \tag{CFT.11}
\]
Norm convergence also gives $\delta_{b_n}(y)\to f(y)$.
Hence $f=\delta_b\in C_0(N)$, so $b\in\mathscr B_p$ and
$(b,f)\in\mathscr G_p$. Metrizability again makes this
sequence argument a proof of closedness.

\emph{The continuous common displacement envelope.}
For the compact quotient family $u\mapsto b_u$ above, define
\[
 a(y)=\max_{u\in P}\delta_{b_u}(y),\qquad
 a^-(y)=\max_{u\in P}\delta_{b_u^{-1}}(y).
                                                               \tag{CFT.12}
\]
The functions $(u,y)\mapsto\delta_{b_u}(y)$ and their inverse
counterparts are jointly continuous. Their pullbacks by
$\id_P\times\pi$ are the continuous functions
$d_Q(h_{b_u}z,z)$ and $d_Q(h_{b_u}^{-1}z,z)$, respectively;
the product covering is open and hence is a quotient map.
Compactness of $P$ then makes both maxima continuous.
Explicitly, at $y_0$ use continuity on product neighborhoods
to bound the change in each value by $\epsilon$, take finitely
many such parameter neighborhoods covering $P$, and intersect
their $y$-neighborhoods. This bounds the change of the maximum
by the same $\epsilon$.

The following properties of this fixed compact quotient
family are equivalent:
\begin{enumerate}
\item $u\mapsto L(b_u)$ is continuous in the original
      closed-disk map-and-inverse topology;
\item $u\mapsto\delta_{b_u}$ is continuous as a map
      $P\to C_0(N)$ with its supremum norm;
\item $a\in C_0(N)$;
\item the family has the uniform tail in (C.92.17).
\end{enumerate}
The first two are equivalent by the proved graph
homeomorphism. Compactness of the image in the second
property and the finite-net tail proof above give the third.
The third and fourth are the literal equality (CFT.12).
For a direct converse from the fourth to the second, choose
one compact tail core with bound $\epsilon/3$. On that core,
joint continuity and compactness give uniform convergence
$\delta_{b_v}\to\delta_{b_u}$ as $v\to u$. Off the core both
nonnegative functions are less than $\epsilon/3$, so their
difference is less than $\epsilon/3$. Their supremum-norm
difference therefore tends to zero. Finally (CFT.6)--(CFT.7)
prove that these properties also give $a^-\in C_0(N)$;
the same argument proves norm continuity of the inverse
displacement family.

For a compact family continuous in the original topology,
there is an explicit strictly positive continuous common
bound in $C_0(N)$. Construct continuous functions $\varphi_n$
equal to one on $K_n$ and supported in $K_{n+1}$ as follows.
When $\operatorname{int}K_{n+1}\ne N$, put
\[
 \varphi_n(y)=
 \frac{\operatorname{dist}_\rho(y,N\setminus\operatorname{int}K_{n+1})}
 {\operatorname{dist}_\rho(y,N\setminus\operatorname{int}K_{n+1})
           +\operatorname{dist}_\rho(y,K_n)}.
                                                               \tag{CFT.13}
\]
The denominator is positive: its two zero sets are the
disjoint closed sets $N\setminus\operatorname{int}K_{n+1}$
and $K_n$. When $\operatorname{int}K_{n+1}=N$, the space $N$
is compact, and take $\varphi_n=1$. In every case
$0\leq\varphi_n\leq1$, $\varphi_n=1$ on $K_n$, and its
support is compact and contained in $K_{n+1}$.
The uniformly convergent series
\[
 \tau(y)=\sum_{n=1}^{\infty}2^{-n}\varphi_n(y),\qquad
 \eta(y)=a(y)+a^-(y)+\tau(y)
                                                               \tag{CFT.14}
\]
has $\tau>0$, because every $y$ lies in some $K_n$.
It belongs to $C_0(N)$: outside $K_{M+1}$ the first $M$
summands vanish and the remaining sum is at most
$\sum_{n>M}2^{-n}$. Thus, for each continuous original
compact family, (CFT.12)--(CFT.14) define an actual continuous
function $\eta\in C_0(N)$ with
\[
 0\leq\delta_{b_u}(y)\leq a(y)<\eta(y),\qquad
 0\leq\delta_{b_u^{-1}}(y)\leq a^-(y)<\eta(y)
                  \quad(u\in P,\ y\in N).
                                                               \tag{CFT.15}
\]

All maps in this comparison retain the unique pointed lifts
and every original deck label. The radial example in
(C.92.18)--(C.92.21), with its explicitly assigned Euclidean
closed-disk metric, has $b_R\to1$ for (CFT.3) and
$\|\delta_{b_R}\|_\infty>2/5$. Thus its graph coordinates
fail to converge to $(1,0)$, exactly as its original lifts
fail to converge to the identity. This verifies the graph
comparison on that previously proved example by the actual
coordinate maps (CFT.9).

The proved output here is the closed graph homeomorphism,
the finite-core and inverse-core estimates, and the explicit
continuous common bound. The noncompact compact-support
deformation in C.88 remains unresolved; an actual deformation
map has not been constructed in this note. C.91 and C.92,
including their stated scopes, remain unchanged.

\end{document}
